\section{用三角板画抛物线}
\markboth{8. 用三角板画抛物线}{8. 用三角板画抛物线}

抛物线也可利用三角板来画，其方法既有趣，又实用。如图 8.1，取一点 $F$，并且作一条不通过 $F$ 的直线 $a$。让三角板的一条直角边 $LM$ 通过定点 $F$，直角顶点 $M$ 放在直线 $a$ 上的任一位置，然后沿另一条直角边 $MN$ 画直线，就得到所要作的抛物线的一条切线。用同样方法再画出若干条直线，最后再画一条与这些直线都相切的光滑曲线，就得到了所要画的抛物线。

\begin{figure}[htbp]
  \centering
  \begin{minipage}[t]{0.47\textwidth}
    \centering
    \begin{tikzpicture}[scale=1.02]
      % 抛物线 y = x^2 / 1.8, 顶点 A=(0,0), 焦点 F=(0, 0.45), 直线 a: y=0
      \coordinate (A) at (0,0);
      \coordinate (F) at (0,0.45);
      % 直线 a
      \draw[dashed, line width=0.65pt] (-1.35,0) -- (1.45,0);
      \node[below] at (1.30,0) {\footnotesize $a$};
      % 抛物线
      \draw[line width=0.9pt, curveblue, domain=-1.52:1.52, samples=70]
        plot (\x, {\x*\x/1.8});
      % 右侧几条切线: 在 a 上取 M=(m,0), 切点 P=(2m, 2m^2/0.9) = (2m, 4m^2/1.8)
      % 斜率为 2m/0.9 = m/0.45, 方程 y = (m/0.45)*(x - m)
      \foreach \m in {0.20, 0.38, 0.55} {
        \draw[line width=0.45pt, black!75] (\m-0.28, {(\m/0.45)*(-0.28)}) -- (2*\m+0.36, {(\m/0.45)*(\m+0.36)});
        \draw[line width=0.4pt, dashed, black!60] (\m,0) -- (F);
      }
      % 三角板: 直角顶点 M=(-0.50, 0), F=(0, 0.45)
      % MF 方向 (0.50, 0.45), 单位化 (0.7433, 0.6690)
      % 垂直方向 (-0.6690, 0.7433)
      \coordinate (M) at (-0.50, 0);
      \coordinate (L) at ($(M) + 1.45*(0.7433, 0.6690)$);
      \coordinate (N) at ($(M) + 1.75*(-0.6690, 0.7433)$);
      \draw[line width=0.7pt, fill=white, fill opacity=0.85] (M) -- (L) -- (N) -- cycle;
      % 三角板内部镂空三角形
      \coordinate (M1) at ($(M) + 0.30*(0.7433, 0.6690) + 0.30*(-0.6690, 0.7433)$);
      \coordinate (L1) at ($(M) + 0.95*(0.7433, 0.6690) + 0.25*(-0.6690, 0.7433)$);
      \coordinate (N1) at ($(M) + 0.25*(0.7433, 0.6690) + 1.05*(-0.6690, 0.7433)$);
      \draw[line width=0.5pt] (M1) -- (L1) -- (N1) -- cycle;

      % 标记点
      \fill (F) circle (1.3pt) node[below right=-1pt] {\footnotesize $F$};
      \fill (A) circle (1.1pt) node[below=1pt] {\footnotesize $A$};
      \fill (M) circle (1.1pt) node[below=1pt] {\footnotesize $M$};
      \node[above right=-1pt] at (L) {\footnotesize $L$};
      \node[left] at (N) {\footnotesize $N$};
    \end{tikzpicture}
    \caption*{图\ 8.1}
  \end{minipage}
  \hfill
  \begin{minipage}[t]{0.47\textwidth}
    \centering
    \begin{tikzpicture}[scale=1.02]
      % 抛物线 y = x^2 / 1.8, 焦点 F=(0, 0.45), 准线 l: y=-0.45, 直线 a: y=0
      \coordinate (F) at (0,0.45);
      \coordinate (A) at (0,0);
      \coordinate (B) at (0,-0.45);
      \coordinate (M) at (-0.56,0);
      \coordinate (G) at (-1.12,-0.45);
      \coordinate (P) at (-1.12,0.6969);
      \coordinate (N) at (-1.52,1.1947);

      % 准线 l 与直线 a, 轴 FB
      \draw[line width=0.65pt] (-1.45,-0.45) -- (1.35,-0.45);
      \node[below] at (1.20,-0.45) {\footnotesize $l$};
      \draw[dashed, line width=0.65pt] (-1.45,0) -- (1.35,0);
      \node[above] at (1.20,0) {\footnotesize $a$};
      \draw[line width=0.55pt] (0,1.25) -- (B);

      % 抛物线
      \draw[line width=0.9pt, curveblue, domain=-1.28:1.25, samples=60]
        plot (\x, {\x*\x/1.8});

      % 辅助线: F-M-G, G-P, P-F, 切线 N-P-M 延长 (斜率 -1.12/0.9 = -1.2444)
      \draw[line width=0.55pt] (F) -- (G);
      \draw[line width=0.55pt] (G) -- (P);
      \draw[line width=0.55pt] (P) -- (F);
      \draw[line width=0.65pt] (N) -- (-0.30,-0.3236);

      % 标记点与直角
      \draw[line width=0.45pt] (G) ++(0,0.12) -- ++(0.12,0) -- ++(0,-0.12);
      \fill (F) circle (1.3pt) node[right=1pt] {\footnotesize $F$};
      \fill (P) circle (1.2pt) node[left=1pt] {\footnotesize $P$};
      \fill (M) circle (1.2pt) node[below=2pt, fill=white, inner sep=0.8pt] {\footnotesize $M(M')$};
      \node[above right=-2pt] at (A) {\footnotesize $A$};
      \node[below] at (B) {\footnotesize $B$};
      \node[below] at (G) {\footnotesize $G$};
      \node[left] at (N) {\footnotesize $N$};
    \end{tikzpicture}
    \caption*{图\ 8.2}
  \end{minipage}
\end{figure}

这种画法，是根据抛物线切线的一个简单性质设计出来的。如图 8.2，设定点 $F$ 到定直线 $l$ 的距离为 $FB$，过 $FB$ 的中点 $A$ 作 $a \parallel l$。在 $a$ 上任取一点 $M$，作 $\angle FMN = 90^\circ$，现证 $MN$ 是抛物线的一条切线。

延长 $FM$ 交 $l$ 于 $G$，则 $FM = MG$。作 $GP \perp l$，交 $MN$ 于 $P$，由于 $\angle FMP = 90^\circ$，所以 $\triangle FMP \cong \triangle GMP$。从而得到 $PG = PF$，即 $P$ 点在以 $F$ 为焦点、$l$ 为准线的抛物线上，并且 $\angle GPM = \angle MPF$。过 $P$ 点作此抛物线的切线 $PM'$，根据切线的性质，$\angle FPM' = \angle GPM'$，可见 $PM$ 与 $PM'$ 重合，$MN$ 就是抛物线在 $P$ 点的切线。显然，$MN$ 是不会平行于 $BF$ 的。
